\documentclass[10pt]{article}
\usepackage[a4paper,margin=14mm]{geometry}
\usepackage[no-math]{fontspec}
\setmainfont{Latin Modern Roman}
\usepackage{amsmath,amssymb,amsthm,amsfonts}
\usepackage{xurl}
\usepackage[unicode,hidelinks]{hyperref}
\setlength\emergencystretch{3em}
% Independently written finite standard notation wrapper. Exact source inventory is checked before build.
\newcommand{\E}{\mathop{\mathbb E}}
\newcommand{\card}[1]{\left\lvert #1\right\rvert}
\newcommand{\set}[1]{{\left\{#1\right\}}}
\newcommand{\sett}[2]{\left\{\left. #1\;\right\vert #2\right\}}
\newcommand{\Prob}[2]{{\Pr_{#1}\left[#2\right]}}
\newcommand{\Expect}[2]{{\mathop{\mathbb E}_{#1}\left[#2\right]}}
\newcommand{\norm}[1]{\lVert #1\rVert}
\newcommand{\power}[1]{\set{0,1}^{#1}}
\newcommand{\defeq}{\stackrel{def}{=}}
\newcommand{\inner}[2]{\langle #1,#2\rangle}
\newcommand{\eps}{\varepsilon}
\renewcommand{\gg}{\gtrsim}
\renewcommand{\ll}{\lesssim}
\DeclareRobustCommand{\etalchar}[1]{\def\ObservedPlus{+}\def\ActualPlus{#1}\ifx\ObservedPlus\ActualPlus\ensuremath{{}^{+}}\else\PackageError{finite-printed-glyph}{Unobserved etalchar argument}{Only the inventoried plus argument is allowed}\fi}
% Ordinary standard amsthm environments with exact authored shared numbering.
\newtheorem{thm}{Theorem}[section]
\newtheorem{lemma}[thm]{Lemma}
\newtheorem{corollary}[thm]{Corollary}
\newtheorem{claim}[thm]{Claim}
\newtheorem{definition}[thm]{Definition}
\newtheorem{fact}[thm]{Fact}
\begin{document}\begin{center}\Large Problem 2451: Sensitivity-moment junta approximation\end{center}\noindent\textbf{Authors:} Ronen Eldan, Guy Kindler, Noam Lifshitz and Dor Minzer\par\noindent\textbf{Work/version:} Isoperimetric Inequalities Made Simpler; Published Discrete Analysis 2025:7, 23 pp.; published 24 July 2025; DOI 10.19086/da.142095; matching arXiv:2204.06686v4 (21 July 2025)\par\noindent\textbf{Published source:} \url{https://arxiv.org/pdf/2204.06686v4}\par\noindent\textbf{Published license:} CC BY 3.0, \url{https://creativecommons.org/licenses/by/3.0/}\par\noindent\textbf{Difficulty:} H2: The proof selects coordinates by noisy influence and decomposes all omitted Fourier mass into dyadic scales. A random-restriction argument separates low and high first-level coefficients, using gradient moments for one part and near-L2 hypercontractivity for the other. A feedback estimate controls the sum of noisy influences by the unknown dyadic mass itself, closing the bound and obtaining the sharp exponential dependence on the moment-to-error ratio.\par\medskip\noindent\textbf{Editorial source boundary:} Theorem 1.5 is the sole selected outcome, with every parameter, the uniform-cube convention and junta definition retained. The complete Fourier/gradient/restriction context, first-level and dyadic tail arguments 3.1-3.3, full Lp argument 3.6 and literal author reuse pointer for 3.7 are included. Lemma 3.7 is stated after the author says "As in Lemma 3.2"; no separate expansion is added, and the exact reused human argument is retained. All genuinely new candidate-coordinate, noisy-influence, soft-truncation, dyadic feedback and rounding steps of Section 4 are included. Claim 4.4 first says "Deferred to Section 4.2"; its complete later proof 1041-1218, including Claim 4.9 and both low/high restriction contributions, is retained. Original background and Corollary 3.11 remain context without extra counts. Named Bonami-Beckner-Gross and O'Donnell hypercontractive/noisy-influence facts stay precisely attributed prerequisites. Printed mathematics is a checked transcription of the CC BY 3.0 published PDF; the matching nonexclusive native archive/preamble/classes and unprinted comments stay private. Standard notation bindings are independently written for the exact finite inventoried calls. Original statement, proof order and mathematical tokens are unchanged.\par\medskip\hrule\medskip\noindent\textbf{Checked published author text follows.}\par
% Checked published text; private-aid LF positions 146--168

\section{Introduction}
\subsection{Isoperimetric Inequalities Over the Hypercube}
Isoperimetric inequalities are fundamental results in mathematics with a myriad of applications. The main topic of this paper is discrete isoperimetric
inequalities, and more specifically isoperimetric inequalities over the hypercube $\{0,1\}^n$. Classical results along
these lines are due to Harper, Bernstein, Lindsey and Hart~\cite{Harper2,Bernstein,Lindsey,Hart}, Harper~\cite{Harper2}, Margulis~\cite{Margulis}, Talagrand~\cite{Tal1,Tal2,Tal4} and Kahn, Kalai and Linial~\cite{KKL}.

All of these theorems discuss various measures of boundaries and relations between them for Boolean functions.
For a function $f\colon\power{n}\to\power{}$ and a point $x\in\power{n}$, we define the \emph{sensitivity} of $f$ at $x$, denoted by $s_f(x)$,
to be the number of coordinates $i\in[n]$ such that $f(x\oplus e_i)\neq f(x)$. Thus, the \emph{vertex boundary} $V_f$ of a function $f$ is defined to be the set
of all $x$'s such that $s_f(x) > 0$, and the \emph{edge boundary} $E_f$ of $f$ is defined to be the set of edges $(x,x\oplus e_i)$ of the hypercube whose
endpoints are assigned different values by $f$. The most basic isoperimetric result related to Boolean functions,
known as Poincar\'{e}'s inequality, asserts that $\frac{\card{E_f}}{2^n}\geq {\sf var}(f)$,
i.e.\ the edge boundary of a function $f$ cannot be too small.\footnote{We remark that the term ``Poincar\'{e}'s inequality'' comes from the theory of functional inequalities
and Markov chains and refers to a much more general result, see for example~\cite[Section 2]{kitsos2014inequalities}. In the case of graphs, it asserts that for a graph
$G$, if the first non-trivial eigenvalue of its Laplacian $L$ is at least $\lambda$, then $\norm{L f}_2\geq \lambda \norm{f-\E[f]}_2$. In our
example, the Laplacian $L$ is $L = I-\frac{1}{n} A$ where $A$ is the adjacency matrix of the Boolean hypercube, and $\lambda$ is $2/n$.}
The edge isoperimetric inequality for the hypercube, due to Harper, Bernstein, Lindsey and Hart~\cite{Harper2,Bernstein,Lindsey,Hart}, is a strengthening
of that statement, asserting that for a fixed value of $\E[f]$, the functions that minimize the edge boundary $\frac{\card{E_f}}{2^n}$
are indicator functions
of initial segments in lexicographical orderings of the hypercube. In particular, this result implies that if $f$ is a non-constant Boolean function, then $\frac{\card{E_f}}{2^n}\geq \E[f]\log\left(\frac{1}{\E[f]}\right)$
(here and throughout, $\log$ means the logarithm function with base $2$). The vertex isoperimetric inequality, due to Harper~\cite{Harper2}, asserts that among all Boolean function
of a given measure, the functions that minimize the size of the vertex boundary $V_f$ are the indicator functions of
initial segments in simplicial orderings of the hypercube.
% Checked published text; private-aid LF positions 172--173

Talagrand considered a different notion of boundary for a function $f$, namely the quantity $\Expect{x}{\sqrt{s_f(x)}}$, which we will refer to as the \emph{Talagrand boundary} of $f$.
Here and throughout, the distribution over $x$ will be uniform over the hypercube $\power{n}$.
% Checked published text; private-aid LF positions 270--304
\setcounter{subsection}{3}\setcounter{thm}{3}
\subsection{New Junta Theorems}
Next, we study connections between Talagrand-like quantities such as $\Expect{x}{\sqrt{s_f(x)}}$, and the notion of juntas.
For a point $x\in \{0,1\}^n$ and a subset $J\subseteq [n]$, we denote by $x_J$ the point in $\{0,1\}^J$ corresponding to
the value of $x$ on coordinates $J$.
\begin{definition}
  A function $f\colon\{0,1\}^n\to\{0,1\}$ is called a $j$-junta if there exists $J\subseteq [n]$ of size at most $j$,
  and $g\colon \{0,1\}^J\to\{0,1\}$ such that $f(x) = g(x_J)$ for all $x\in\{0,1\}^n$.

  We say a function $f\colon\{0,1\}^n\to\{0,1\}$ is $\eps$-close to a $j$-junta if there is a $j$-junta $g\colon\{0,1\}^n\to\{0,1\}$
  such that $\Prob{x}{f(x)\neq g(x)}\leq \eps$.
\end{definition}
On the one hand, it is possible that the quantity $\Expect{x}{\sqrt{s_f(x)}}$ is constant while the function $f$ is very far from being a junta,
as can be seen by taking $f$ to be the majority function. On the other hand, if we look at the expected sensitivity -- that is
$\Expect{x}{s_f(x)}$ -- a well known result of Friedgut~\cite{Friedgut} asserts that if $\Expect{x}{s_f(x)} = O(1)$, then $f$
is close to junta. More precisely, Friedgut's result asserts that if $\Expect{x}{s_f(x)} \leq A$, then
for all $\eps>0$, $f$ is $\eps$-close to a $2^{O(A/\eps)}$-junta. In light of this contrast, it is interesting to ask
what happens if an intermediate moment of the sensitivity, say the $p$th moment for $p\in (1/2,1)$, is constant.

We prove a strengthening of Friedgut's junta theorem~\cite{Friedgut} addressing this case, morally showing that for any $p>1/2$, bounded $p$-moment of
the sensitivity implies closeness to junta. More precisely:
\begin{thm}\label{thm:super_fried}
  For all $\eta>0$ there is $C = C(\eta)>0$ such that  the following holds.
  For all $\eps>0$ and $A>0$, and $1/2 + \eta\leq p\leq 1$, if $f\colon\{0,1\}^n\to\{0,1\}$ is a function such that
  $\Expect{x}{s_f(x)^{p}}\leq A$, then $f$ is $\eps$-close to a $J$-junta, where $J = 2^{C\left(A/\eps\right)^{1/p}}$.
\end{thm}
Theorem~\ref{thm:super_fried} is tight, and can be seen as a generalization of Friedgut's Junta Theorem (corresponding to the case that
$p=1$) that is somewhat close in spirit to Bourgain's Junta Theorem. Indeed, for all $p\geq 1/2$ the tribes function shows
this to be tight. Here, the tribes function is the function in which we take a partition $P_1\cup\ldots\cup P_m$ of $[n]$
and define ${\sf Tribes}(x) = 1$ if there is $i$ such that $x_{P_{i}} = \vec{1}$, and ${\sf Tribes}(x) = 0$ otherwise.
We take $P_i$ of equal size and the probability that ${\sf Tribes}(x)=1$ is roughly $1/2$; the correct choice
of parameters is $\card{P_i} = \log n - \log\log n + c$
and $m =  \frac{n}{\log n - \log\log n + c}$ where $c = O(1)$.
For this function one has $\Expect{x}{s_{{\sf Tribes}}(x)^{p}} = \Theta((\log n)^p)$. Indeed, as $s_{{\sf Tribes}}(x) = \Theta(\log n)$ with
constant probability we have the lower bound, and the upper bound follows by Jensen's inequality and the fact that the total influence of ${\sf Tribes}$
is $\Theta(\log n)$.
% Checked published text; private-aid LF positions 415--507

\section{Preliminaries}
\paragraph{Notation.} We will use big $O$ notations throughout the paper. Sometimes, it will be easier for us to use
$\ll$ and $\gg$ notations; when we say that $X\ll Y$ we mean that there is an absolute constant $C>0$ such that $X\leq C Y$.
Analogously, when we say that $X\gg Y$ we mean that there is an absolute constant $C>0$ such that $X\geq C Y$.

\subsection{Fourier Analysis over the Hypercube}
We consider the space of real-valued functions $f:\power{n} \to {\mathbb R}$,
equipped with the inner product $\inner{f}{g} = \Expect{x\in_R\power{n}}{f(x)g(x)}$.
It is well-known that the collection of functions $\set{\chi_S}_{S\subseteq[n]}$ defined by
$\chi_S(x) = \prod\limits_{i\in S}{(-1)^{x_i}}$ forms an orthonormal basis, so any function $f$
may be written as $f(x) = \sum\limits_{S\subseteq[n]}{\widehat{f}(S)\chi_S(x)}$ in a unique way,
where
$\widehat{f}(S) = \inner{f}{\chi_S}$.
We also consider $L_p$ norms for $p \geq 1$,
that are defined as $\norm{f}_p= \left(\Expect{x}{\card{f(x)}^p}\right)^{1/p}$.

For $z\in\power{n}$, we denote by $z_{-i}\in\power{n-1}$ the vector resulting from dropping
the $i$th coordinate of $z$. We denote by $(x_i = 1, x_{-i} = z_{-i})$ the $n$-bit vector
which is $1$ on the $i$th coordinate and is equal to $z$ on any other coordinate.
\begin{definition}
  The derivative of $f$ along direction $i\in[n]$
  is $\partial_i f\colon\power{n}\to\mathbb{R}$ defined by
  $\partial_i f(z) = f(x_i = 0, x_{-i} = z_{-i}) - f(x_i = 1, x_{-i} = z_{-i})$.
  The gradient of $f$, $\nabla f\colon\power{n}\to\mathbb{R}^n$, is defined
  as $\nabla f(z) = (\partial_1 f(z),\ldots,\partial_n f(z))$.
\end{definition}
Note that for a Boolean function $f$, for each $x\in\power{n}$, $\partial_i f(x) \neq 0$ if and only if the coordinate $i$
is sensitive on $x$, and in this case its value is either $1$ or $-1$. Thus, $\norm{\nabla f(x)}_2 = \sqrt{s_f(x)}$,
and it will be often convenient for us to use the $\ell_2$-norm of the gradient of $f$ in place of $\sqrt{s_f(x)}$.

\begin{fact}\label{fact:singeton}
  Let $f\colon\power{n}\to\mathbb{R}$, and let $i\in[n]$.
  Then $\partial_i f(x) =
  2\sum\limits_{S\ni i}{\widehat{f}(S)\chi_{S\setminus \set{i}}(x)}$.
  In particular, we have that $\Expect{x}{\partial_i f(x)} = 2\widehat{f}(\set{i})$.
\end{fact}


Next, we define the level $d$ \emph{Fourier weight} of a function $f$, the \emph{approximate level $d$ weight} of a function $f$
and the \emph{Fourier weight of $f$ up to/ above level $d$}.
\begin{definition}
  The level $d$ Fourier weight of $f$ is defined to be $W_{=d}[f] = \sum\limits_{\card{S} = d}{\widehat{f}(S)^2}$. The
  approximate level $d$ weight of a function $f$ is defined to be $W_{\approx d}[f] = \sum\limits_{d\leq j<2d}{W_{=j}[f]}$.
  The weight of $f$ up to level $d$ is defined as $W_{\leq d}[f] = \sum\limits_{\card{S}\leq d}{\widehat{f}(S)^2}$,
  and the weight of $f$ above level $d$ is defined as $W_{> d}[f] = \sum\limits_{\card{S}> d}{\widehat{f}(S)^2}$.
\end{definition}
For a function $f\colon\power{n}\to\mathbb{R}$, we also define the \emph{level at most $d$ part} of $f$, denoted by $f^{\leq d}$,
as $f^{\leq d}(x) = \sum\limits_{\card{S}\leq d} \widehat{f}(S)\chi_S(x)$. We note that by Parseval's equality we have that
$\norm{f^{\leq d}}_2^2 = W_{\leq d}[f]$.
\subsection{Random Restrictions}
For a function $f\colon\{0,1\}^n\to\mathbb{R}$, a set of coordinates $J\subseteq[n]$ and $z\in\power{\overline{J}}$, we define the
\emph{restriction} of
$f$ according to $(J,z)$ as the function $f_{\overline{J}\rightarrow z}\colon \{0,1\}^{J}\to\mathbb{R}$ defined by
\[
f_{\overline{J}\rightarrow z}(y) = f(x_{\overline{J}} = z, x_{J} = y).
\]
We refer to the coordinates of $J$ as ``alive'', and refer to the coordinates of $\overline{J}$ as fixed.
We will often be interested in random restrictions of a function $f$. By that, we mean that the set $J$ is chosen randomly by including in it each
$i\in[n]$ independently with some probability (which is specified), and $z\in\power{\overline{J}}$ is chosen uniformly.
We have the following fact, which establishes a relation between
the approximate level $d$ weight of $f$ and the level $1$ weight of a random restriction of $f$ in which each $i\in [n]$ is included in $J$ with probability
$\frac{1}{2d}$ (see~\cite[Fact 4.1]{kindler2018gaussian} and~\cite[Section 3.1]{KKKMS}).
\begin{fact}\label{fact:restriction}
  Let $f\colon\power{n}\to\mathbb{R}$ and $d\in\mathbb{N}$, let $(J,z)$ be a random restriction
  where each $j\in [n]$ is included in $J$ with probability $\frac{1}{2d}$, and $z\in\power{\bar{J}}$ is chosen uniformly. Then $\Expect{J,z}{W_{=1}[f_{\bar{J}\rightarrow z}]}\gg W_{\approx d}[f]$.
\end{fact}
\begin{proof}
  For a fixed $J$ and $T\subseteq J$, we have that
  $\widehat{f_{\bar{J}\rightarrow z}}(T) = \sum\limits_{S\subseteq \bar{J}}\widehat{f}(T\cup S)\chi_S(z)$. Thus,
  \[
  \Expect{J,z}{W_{=1}[f_{\bar{J}\rightarrow z}]}
  =\Expect{J}{\Expect{z}{\sum\limits_{j\in J}\left(\sum\limits_{S\subseteq\bar{J}}\widehat{f}(S\cup \set{j})\chi_S(z)\right)^2}}
  =\Expect{J}{\sum\limits_{j\in J}\sum\limits_{S\subseteq\bar{J}}\widehat{f}(S\cup \set{j})^2},
  \]
  where the last transition follows by pushing the expectation over $z$ inside and using Parseval's equality. It follows that
  \[
  \Expect{J,z}{W_{=1}[f_{\bar{J}\rightarrow z}]} =
  \Expect{J}{\sum\limits_{T}\widehat{f}(T)^21_{\card{T\cap J} = 1}}
  =\sum\limits_{T}\widehat{f}(T)^2\Prob{J}{\card{T\cap J}=1},
  \]
  and as $\Prob{J}{\card{T\cap J}=1}\gg 1$ for all $T$ whose size is between $d$ and $2d$ and the rest of the summands
  are all non-negative, the proof is concluded.
\end{proof}

\subsection{Bonami's Lemma}
The \emph{degree} of a function $f\colon \power{n}\to\mathbb{R}$ is defined as
${\sf deg}(f) = \max\sett{\card{S}}{\widehat{f}(S)\neq 0}$.
We will need the following consequence of
the Bonami-Beckner-Gross hypercontractive inequality~\cite{Bec75,Bonami,grossLogSob} (see also~\cite{ODonnell}), showing that the $4$-norm of a low-degree function $f$ is comparable to its $2$-norm.
\begin{thm}\label{thm:hypercontractivity}
  Let $f\colon\power{n}\to\mathbb{R}$ be a function of degree at most $d$. Then
  $\norm{f}_4\leq \sqrt{3}^d\norm{f}_2$.
\end{thm}
% Checked published text; private-aid LF positions 511--511

\section{Isoperimetric Inequalities on the Hypercube}\label{sec:classic_iso}
% Checked published text; private-aid LF positions 514--572
\setcounter{equation}{1}
\begin{lemma}\label{lemma:tal_lvl1}
  Let $f\colon\power{n}\to\mathbb{R}$. Then
  $\Expect{}{\norm{\nabla f(x)}_2 }\geq 2\sqrt{W_{=1}[f]}$.
\end{lemma}
\begin{proof}
By the triangle inequality and Fact~\ref{fact:singeton} we have
  \[
  \Expect{x}{\norm{\nabla f(x)}_2 }
  \geq \norm{\Expect{x}{\nabla f(x)}}_2
  = \norm{2(\widehat{f}(\set{1}),\ldots,\widehat{f}(\set{n}))}_2
  = 2\sqrt{\sum\limits_{i=1}^{n}\widehat{f}(\set{i})^2}
  = 2\sqrt{W_{=1}[f]}.
  \qedhere
 \]
\end{proof}

\begin{lemma}\label{lemma:tal_lvld_exact}
  Let $d\in\mathbb{N}$, and $f\colon\power{n}\to\power{}$. Then
  $\Expect{}{\norm{\nabla f(x)}_2 }\gg \sqrt{d}W_{\approx d}[f]$.
\end{lemma}
\begin{proof}
  Let $(J,z)$ be a random restriction for $f$, where each $j\in[n]$ is included in $J$ independently with probability
  $\frac{1}{2d}$, and $z\in\power{\bar{J}}$ is sampled uniformly. Fix $x\in \power{n}$, and note that
  \[
    \Expect{J}{\norm{\nabla f_{\bar{J}\rightarrow x_{\bar{J}}}(x_J)}_2}
    \leq \sqrt{\Expect{J}{\norm{\nabla f_{\bar{J}\rightarrow x_{\bar{J}}}(x_J)}^2_2}}
    =\sqrt{\Expect{J}{\sum\limits_{j}{\card{\partial_j f(x)}^2 1_{j\in J}}}}
    =\frac{\norm{\nabla f(x)}_2}{\sqrt{2d}}.
  \]
  Therefore,
  \begin{equation}\label{eq1}
  \frac{1}{\sqrt{2d}}\Expect{}{\norm{\nabla f(x)}_2}\geq \Expect{J,z}{\Expect{x\in\power{J}}{\norm{\nabla f_{\bar{J}\rightarrow z}(x)}_2}}.
  \end{equation}
  On the other hand, for each $J$ and $z\in\power{\bar{J}}$ we may apply Lemma~\ref{lemma:tal_lvl1} on $f_{\bar{J}\rightarrow z}$ to get that
  \[
  \Expect{x\in\power{J}}{\norm{\nabla f_{\bar{J}\rightarrow z}(x)}_2}\gg \sqrt{W_{=1}[f_{\bar{J}\rightarrow z}]}\geq W_{=1}[f_{\bar{J}\rightarrow z}].
  \]
  Taking expectation over $J$ and $z$ we get that
  \[
  \Expect{J,z}{\Expect{x\in\power{J}}{\norm{\nabla f_{\bar{J}\rightarrow z}(x)}_2}}
  \gg \Expect{J,z}{W_{=1}[f_{\bar{J}\rightarrow z}]}
  \gg W_{\approx d}[f],
  \]
  where the last inequality is by Fact~\ref{fact:restriction}. Combining this with equation~\eqref{eq1} finishes the proof.
\end{proof}


\begin{lemma}\label{lemma:tal_lvld_above}
  Let $d\in\mathbb{N}$, and $f\colon\power{n}\to\power{}$. Then
  $\Expect{}{\norm{\nabla f(x)}_2 }\gg \sqrt{d}W_{\geq d}[f]$.
\end{lemma}
\begin{proof}
  Applying Lemma~\ref{lemma:tal_lvld_exact} with $d 2^{j}$ in place of $d$ for $j=0,1,2,\ldots$, we get that
  \[
    2^{-j/2}\Expect{}{\norm{\nabla f(x)}_2 }
    \gg \sqrt{d}W_{\approx 2^{j} d}[f],
  \]
  and summing over $j$ yields the result.
\end{proof}
% Checked published text; private-aid LF positions 654--675
\setcounter{subsection}{2}\setcounter{thm}{5}
\subsection{Extensions}
\subsubsection{Other $L_p$-norms}
Lemma~\ref{lemma:tal_lvl1} applies not only in the case of $L_2$-norms, but rather for any $L_p$ norm:
\begin{lemma}\label{lemma:tal_lvl1_L_p}
  Let $1\leq p\leq 2$ and $f\colon\power{n}\to\power{}$. Then
  $\Expect{}{\norm{\nabla f(x)}_p }\geq W_{=1}[f]^{1/p}$.
\end{lemma}
\begin{proof}
By Jensen's inequality and Lemma~\ref{lemma:tal_lvl1} we have
  \[
  \Expect{x}{\norm{\nabla f(x)}_p }
  =\Expect{x}{\norm{\nabla f(x)}_2^{2/p}}
  \geq \Expect{x}{\norm{\nabla f(x)}_2}^{2/p}
  \geq W_{=1}[f]^{1/p}.
  \qedhere
  \]
\end{proof}
As in Lemma~\ref{lemma:tal_lvld_exact}, the previous lemma quickly leads to the following conclusion.
\begin{lemma}\label{lemma:tal_lvld_above_alpha}
  Let $d\in\mathbb{N}$, $\alpha\in[1/2, 1]$ and $f\colon\power{n}\to\power{}$.
  Then $\Expect{}{s_f(x)^{\alpha}}\gg d^{\alpha} W_{\geq d}[f]$.
\end{lemma}
% Checked published text; private-aid LF positions 701--748
\setcounter{subsubsection}{1}\setcounter{thm}{9}
\subsubsection{Talagrand Boundary and Noise Sensitivity}
We next demonstrate the connection between having small Talagrand-type boundary and being noise stable, a notion that we define next.

For a parameter $\rho\in [0,1]$ and a point $x\in \power{n}$, the distribution of $\rho$-correlated inputs with $x$, denoted as $y\sim_{\rho} x$,
is defined by the following randomized process. For each $i\in[n]$ independently, set $y_i = x_i$ with probability $\rho$, and otherwise sample $y_i\in\{0,1\}$ uniformly.
The operator $T_{\rho}\colon L_2(\power{n}) \to L_2(\power{n})$, known as the noise operator with noise rate
$1-\rho$, is defined as
\[
T_{\rho} f(x) = \Expect{y\sim_{\rho} x}{f(y)}.
\]
\begin{definition}
  For $\eps>0$, the noise stability of $f$ with noise rate $\eps$ is defined as ${\sf Stab}_{1-\eps}(f) = \inner{f}{T_{1-\eps} f}$.
\end{definition}

The relation established in Lemma~\ref{lemma:tal_lvld_above_alpha} between the Fourier weight of $f$ and $\Expect{}{s_f(x)^{\alpha}}$
for $1/2\leq\alpha\leq 1$ implies a relation between the latter quantity and noise stability, as follows:
\begin{corollary}\label{corr:noise_stable}
  There exists an absolute constant $C>0$, such that for any $\delta\in[0,1/2]$ and for any function $f\colon\power{n}\to\power{}$ satisfying
  $\Expect{}{s_f(x)^{\frac{1}{2}+\delta}}\leq A$, we have that ${\sf Stab}_{1-\eps}(f)\geq 1 - C\cdot A\cdot \eps^{\frac{1}{2} + \delta}\log(1/\eps)$.
\end{corollary}
\begin{proof}
  \[
  1-{\sf Stab}_{1-\eps}(f) = \Prob{x,y\text{ $(1-\eps)$-correlated}}{f(x)\neq f(y)}
  =2\inner{1-f}{T_{1-\eps} f}
  =2\sum\limits_{k=0}^n{(1-(1-\eps)^{k})W_{=k}[f]}.
  \]

  Set $d = 1/\eps$; we split the sum into $k< d$ and $k \geq  d$, and upper bound the contribution of each part separately.
  For $k\geq d$, using Lemma~\ref{lemma:tal_lvld_above_alpha} we have
  \[
  2\sum\limits_{k=d}^n{(1-(1-\eps)^{k})W_{=k}[f]}
  \leq 2W_{\geq d}[f]\ll A\cdot d^{-(1/2+\delta)} = A\eps^{1/2+\delta}.
  \]
  For $k\leq d$, we write
  \[
  \sum\limits_{k=0}^d{(1-(1-\eps)^{k})W_{=k}[f]}
  \leq\sum\limits_{k=1}^d{k\eps W_{=k}[f]}
  \leq \eps \sum\limits_{k=1}^{d} W_{\geq k}[f].
  \]
  Using Lemma~\ref{lemma:tal_lvld_above_alpha} we get $W_{\geq k}[f]\ll A\cdot k^{-(1/2+\delta)}$, and so
  \[
  \sum\limits_{k=0}^d{(1-(1-\eps)^{k})W_{=k}[f]}
  \ll A\cdot \eps \sum\limits_{k=1}^{d} \frac{1}{k^{1/2+\delta}}
  \ll A \cdot \eps d^{1/2-\delta} \sum\limits_{k=1}^{d} \frac{1}{k}
  \ll A \cdot \eps d^{1/2-\delta} \log d.
  \qedhere
  \]
\end{proof}
% Checked published text; private-aid LF positions 826--1218
\setcounter{equation}{2}
\section{Improved Junta Theorems}\label{sec:super_fried}
\subsection{Proof of Theorem~\ref{thm:super_fried}}
In this section, we prove Theorem~\ref{thm:super_fried}. As explained in the introduction, if one does not care about getting tight dependency
between the junta size $J$ and the parameter $A$, then one may use Corollary~\ref{corr:noise_stable} and Bourgain's noise sensitivity theorem~\cite{Bourgain,kindler2018gaussian}
and get a junta theorem with slightly worse parameters than those stated in Theorem~\ref{thm:super_fried}. The main goal of this section
is to get an optimal dependency of $J$ on $A$.

For the proof, we will need the notion of noisy influences of a function $f$, and we first generalize the notion of influences to real valued functions.
\begin{definition}
  Let $f\colon\power{n}\to\mathbb{R}$ be a function, and let $i\in [n]$ be a coordinate.
  We define the influence of $i$ on $f$ as $I_i[f] = \norm{\partial_i f}_2^2$, where
  we recall that $\partial_i f$ is the discrete derivative of $f$ along direction $i$.
\end{definition}
\begin{definition}
  The $\rho$-noisy influence of a coordinate $i$ for $f\colon\{0,1\}^n\to\mathbb{R}$ is defined to be
  $I_i^{(\rho)}[f] = I_i[T_{\rho} f]$.
\end{definition}

\begin{fact}\label{fact:sum_of_noisy}[\cite{ODonnell}]
  For all $f\colon\{0,1\}^n\to\{0,1\}$ we have that $\sum\limits_{i=1}^{n} I_i^{(\rho)}[f] \ll \frac{1}{1-\rho}$.
  \qed
\end{fact}


Let $C_1, C_2$ be two absolute constants, and we think of $C_2$ as much larger than $C_1$. Given $f\colon\{0,1\}^n\to\{0,1\}$ with
$\Expect{x}{s_f(x)^p}\leq A$, we denote $d = C_1\left(\frac{A}{\eps}\right)^{1/p}$, $\delta = 2^{-\frac{C_2}{2p-1} d}$.
Define the candidate set of junta coordinates as
\[
\mathcal{J} = \sett{i\in [n]}{I_i[T_{1-\frac{1}{3d}}[f]]\geq \delta},
\]
i.e.~$\mathcal{J}$ is the set of coordinates whose noisy influence at $f$ is somewhat large. Our task is to show that $f$ is $\eps$-close to a $\mathcal{J}$ junta,
i.e.~that
\begin{equation}\label{eq6}
\sum\limits_{S: S\not\subseteq \mathcal{J}}\widehat{f}(S)^2\leq \eps,
\end{equation}
and the rest of the proof is devoted towards this goal. We remark that the junta $\mathcal{J}$ is
similar in spirit to the junta in Bourgain’s original proof~\cite{Bourgain}, which is easier to see in~\cite[Section 4]{kindler2018gaussian} and
in particular in~\cite[Theorem 4.20]{kindler2018gaussian}. Indeed, both here and in~\cite[Section 4]{kindler2018gaussian}, the junta is
taken to be the set of coordinates with significant noisy influence. There are many other similarities between the argument
presented therein and the argument here. First, in both cases one considers a dyadic sort of partitioning of the Fourier coefficients:
in~\cite[Section 4]{kindler2018gaussian} it is achieved via considering the difference in the noise sensitivity of $f$ at different
noise rates (see~\cite[Theorem 4.16]{kindler2018gaussian}). In the current proof, we perform explicit dyadic partitioning and also
accompany it by applications of suitable noise operators (the operators $S_k$, $S_k'$ and $S_k''$ below). Second, both proofs also use some relation between the noise sensitivity / moments
of the sensitivity and the level $1$ weight of the function $f$ (and its restrictions). While the technical execution of the proofs is
different, it is hard to pinpoint the high level difference between them.
For example, the ``harsh'' degree truncations
in~\cite{kindler2018gaussian} have to be replaced by the ``softer'' type of degree truncations that are achieved by the operators $S_k$, $S_k'$ and
$S_k''$ (see also the parameter $\tilde{W}_{k,d}$) to make the argument go through. Additionally, towards the end of the argument we are able to bound the sum of noisy influences in a non-trivial way, which ultimately allows us to choose $\delta = 2^{O_p(d)}$ as opposed to
$\delta = 2^{O_{p}(d\log d)}$ (which would have resulted in slightly worse bounds as in~\cite[Theorem 4.20]{kindler2018gaussian}); see Claim~\ref{claim:bound_low_deg_inf}.


Towards proving~\eqref{eq6}, we partition the contribution of different scales of
$\card{S\cap \mathcal{J}}$ to the left hand side, and for each $k=1,\ldots,\log d$ define
\[
W_{k,d}[f] = \sum\limits_{\substack{S: \card{S}\leq d,\\ 2^{k-1}\leq \card{S\cap \bar{\mathcal{J}}}\leq 2^k}}\widehat{f}(S)^2.
\]
Thus, we may write
\begin{equation}\label{eq:super_frid1}
  \sum\limits_{S: S\not\subseteq \mathcal{J}}\widehat{f}(S)^2
  \leq
  W_{>d}[f]
  +
  \sum\limits_{k=1}^{\log d} W_{k,d}[f].
\end{equation}
The rest of the proof is devoted to bounding $W_{>d}[f]$ and $W_{k,d}[f]$ for each $k$.
From Lemma~\ref{lemma:tal_lvld_above_alpha} it follows that
\begin{equation}\label{eq:super_frid}
  W_{>d}[f]\leq O\left(\frac{A}{d^p}\right) \leq \frac{\eps}{16}
\end{equation}
for sufficiently large $C_1$. Next, we bound $W_{k,d}[f]$. Fix $k$, and consider the operator $S_k = T_{1-1/d}^{\mathcal{J}}T_{1-1/2^{k}}^{\bar{\mathcal{J}}}$,
i.e. the noise operator that applies $1/2^k$ noise on coordinates outside $\mathcal{J}$, and $1/d$ noise on coordinates of $\mathcal{J}$. Note that
\[
W_{k,d}[f] \ll W_{k,d}[S_k f],
\]
so it suffices to bound $W_{k,d}[S_k f]$. We partition $W_{k,d}[S_k f]$ according to contribution of each $j\in\bar{\mathcal{J}}$, getting:
\begin{align*}
W_{k,d}[S_k f]
=\sum\limits_{\substack{S: \card{S}\leq d,\\ 2^{k-1}\leq \card{S\cap \bar{\mathcal{J}}}\leq 2^k}} \widehat{S_k f}(S)^2
\leq \frac{1}{2^{k-1}} \sum\limits_{j\not\in\mathcal{J}}\sum\limits_{\substack{S\ni j: \card{S}\leq d,\\ 2^{k-1}\leq \card{S\cap \bar{\mathcal{J}}}\leq 2^k}} \widehat{S_k f}(S)^2
&\ll \frac{1}{2^{k-1}} \sum\limits_{j\not\in\mathcal{J}}\norm{S_k \partial_j f}_2^2\\
&\defeq \tilde{W}_{k,d}[f].
\end{align*}
Thus, we have that $W_{k,d}[f]\ll \tilde{W}_{k,d}[f]$, and the rest of the proof is devoted to upper bounding the right hand side.
Define the operators $S_{k}'$ and $S_k''$ as
\[
S_k' = T_{1-\frac{1}{3d}}^{\mathcal{J}}T_{1-\frac{1}{2^{k}}+\frac{2}{3d}}^{\bar{\mathcal{J}}},
\qquad
S_k'' = T_{1-\frac{2}{3d}}^{\mathcal{J}}T_{1-\frac{1}{2^{k}}+\frac{1}{3d}}^{\bar{\mathcal{J}}}.
\]

We establish the following claim using random restrictions and hypercontractivity:
\begin{claim}\label{claim:super_fried_main}
  For all $\tau\in (0,1)$, we have that
  \[
  2^k \tilde{W}_{k,d}[f]\leq
  C\cdot d^{1-p}\tau^{2p-1} \Expect{x}{s_f(x)^p} + C\cdot\tau^{-\frac{1}{100 d}}\sum\limits_{j\in\bar{\mathcal{J}}} \norm{S_k' \partial_j f}_2^{2+\frac{1}{100d}},
  \]
  where $C$ is an absolute constant.
\end{claim}
\begin{proof}
  Deferred to Section~\ref{sec:missing_proofs}.
\end{proof}

Given Claim~\ref{claim:super_fried_main}, one may pull out $\norm{S_k' \partial_j f}_2^{\frac{1}{100 d}}\leq \delta^{\frac{1}{200d}}$ outside
the sum, and bound the sum on the rest of the noisy influences by $O(d)$ to get a bound of the form
$2^k \tilde{W}_{k,d}[f]\ll d^{1-p}\tau^{2p-1} \Expect{x}{s_f(x)^p} + d\tau^{-\frac{1}{100 d}}\delta^{\frac{1}{200d}}$.
This idea is good enough to establish a junta theorem, yet it gives worse parameters.\footnote{One has to take
$\delta = 2^{-O(d\log d)}$, as opposed to $\delta = 2^{-O(d)}$ as we have taken. We remark that if one only cares
about getting a larger junta of size $2^{O(d\log d)}$, the proof greatly simplifies.} Instead, to get the tight result with respect
to the junta size, we will need a more careful bound on the sum of noisy influences outside $\mathcal{J}$. Specifically,
we use an improved bound of the sum of noisy influences that relates them to $W_{\ell,d}[f]$.
To state this bound we define
\[
\eps_{\ell,d} = \sum\limits_{\substack{\card{S}>d\\ 2^{\ell-1}\leq \card{S\cap \bar{\mathcal{J}}}\leq 2^{\ell}}} \widehat{f}(S)^2.
\]
 \begin{claim}\label{claim:bound_low_deg_inf}
    $\frac{1}{2^k}\sum\limits_{j\in\bar{\mathcal{J}}} \norm{S_k' \partial_j f}_2^{2}\ll \sum\limits_{\ell=1}^{\log n} 2^{-\card{\ell-k}}(W_{\ell,d}[f]+\eps_{\ell, d})$.
  \end{claim}
  \begin{proof}
    By definition,
    \[
    \frac{1}{2^k}\sum\limits_{j\in\bar{\mathcal{J}}} \norm{S_k' \partial_j f}_2^{2}
    =4\sum\limits_{S} \frac{\card{S\cap \bar{\mathcal{J}}}}{2^k} \left(1-\frac{1}{2^k}+\frac{2}{3d}\right)^{\card{S\cap \bar{\mathcal{J}}}-1}\left(1-\frac{1}{3d}\right)^{\card{S\cap \mathcal{J}}}\widehat{f}(S)^2.
    \]
    Partitioning the last sum according to $\card{S\cap \bar{\mathcal{J}}}$ and dyadically partitioning so that $2^{\ell-1}\leq \card{S\cap \bar{\mathcal{J}}} < 2^{\ell}$, we have that the last
    expression is at most
    \[
    \ll
    \sum\limits_{\ell=1}^{\log n} 2^{\ell-k} \left(1-\frac{1}{2^k}+\frac{2}{3d}\right)^{2^{\ell-1}}(W_{\ell,d}[f]+\eps_{\ell, d})
    \ll
    \sum\limits_{\ell=1}^{\log n} 2^{-\card{\ell-k}}(W_{\ell,d}[f]+\eps_{\ell, d}).\qedhere
    \]
  \end{proof}

    Combining Claims~\ref{claim:super_fried_main} and~\ref{claim:bound_low_deg_inf} we immediately get the following corollary:
\begin{corollary}\label{corr:bound_low_deg_inf}
  There exists an absolute constant $C>0$ such that for all $\tau>0$,
    \[
    \tilde{W}_{k,d}[f]\leq C\cdot 2^{-k} d^{1-p}\tau^{2p-1} \Expect{x}{s_f(x)^p}
    +C\cdot \tau^{-\frac{1}{100 d}}\delta^{\frac{1}{200 d}}
    \left(\sum\limits_{\ell=1}^{\log d} 2^{-\card{\ell-k}}W_{\ell,d}[f]+\sum\limits_{\ell=1}^{\log n} 2^{-\card{\ell-k}}\eps_{\ell, d}\right).
    \]
\end{corollary}


  We may now give an upper bound on $\sum\limits_{k=1}^{\log d} W_{k,d}[f]$.
  \begin{claim}\label{claim:bd_outside_junta}
  For sufficiently large absolute constant $C_1$ and sufficiently large $C_2$ in comparison to $C_1$, we have that
    $\sum\limits_{k=1}^{\log d} W_{k,d}[f]\leq \frac{\eps}{16}$.
  \end{claim}
  \begin{proof}
    Let $\tau>0$ to be determined. As $W_{k,d}[f]\ll \tilde{W}_{k,d}[f]$, by Corollary~\ref{corr:bound_low_deg_inf} we have
    \[
    \sum\limits_{k=1}^{\log d} W_{k,d}[f]
    \ll
    \sum\limits_{k=1}^{\log d}
    \left(2^{-k}d^{1-p}\tau^{2p-1} A
    +\tau^{-\frac{1}{100 d}}\delta^{\frac{1}{200 d}}
    \left(\sum\limits_{\ell=1}^{\log d} 2^{-\card{\ell-k}}W_{\ell,d}[f]+\sum\limits_{\ell=1}^{\log n} 2^{-\card{\ell-k}}\eps_{\ell, d}\right)
    \right).
    \]
    The first sum is clearly $\ll d^{1-p}\tau^{2p-1} A$. The second sum is at most
    \[
    \tau^{-\frac{1}{100 d}}\delta^{\frac{1}{200 d}}
    \sum\limits_{\ell=1}^{\log d}W_{\ell,d}[f] \sum\limits_{k} 2^{-\card{\ell-k}}
    \ll \tau^{-\frac{1}{100 d}}\delta^{\frac{1}{200 d}}\sum\limits_{\ell=1}^{\log d} W_{\ell,d}[f].
    \]
    Finally, the third sum is at most
    \[
    \tau^{-\frac{1}{100 d}}\delta^{\frac{1}{200 d}}
    \sum\limits_{\ell=1}^{\log n} \eps_{\ell, d}
    \sum\limits_{k=1}^{\log d} 2^{-\card{\ell-k}}
    \ll
    \tau^{-\frac{1}{100 d}}\delta^{\frac{1}{200 d}}
    \sum\limits_{\ell=1}^{\log n} \eps_{\ell, d}
    \ll
    \tau^{-\frac{1}{100 d}}\delta^{\frac{1}{200 d}}
    W_{>d}[f]
    \ll \tau^{-\frac{1}{100 d}}\delta^{\frac{1}{200 d}}\eps,
    \]
    where we used~\eqref{eq:super_frid}.
    We thus get that
    \[
    \sum\limits_{k=1}^{\log d} W_{k,d}[f]
    \leq
    C\cdot d^{1-p}\tau^{2p-1} A
    +C\cdot \tau^{-\frac{1}{100 d}}\delta^{\frac{1}{200 d}}\sum\limits_{k=1}^{\log d} W_{k,d}[f]
    +C\cdot \tau^{-\frac{1}{100 d}}\delta^{\frac{1}{200 d}}\eps
    \]
    for some absolute constant $C>0$. We choose $\tau = \delta^{1/4}$ and then $C_2$ large enough so that
    $C\tau^{-\frac{1}{100 d}}\delta^{\frac{1}{200 d}}\leq \frac{1}{2}$, to get that
    $\sum\limits_{k=1}^{\log d} W_{k,d}[f]\leq 2C\cdot d^{1-p}\delta^{(2p-1)/4} A + 2C\delta^{\frac{1}{400 d}}\eps$,
    which is at most $\frac{\eps}{16}$ for large enough $C_2$.
  \end{proof}

  We can now prove Theorem~\ref{thm:super_fried}.
  \begin{proof}[Proof of Theorem~\ref{thm:super_fried}]
    Combining~\eqref{eq:super_frid1},~\eqref{eq:super_frid} and Claim~\ref{claim:bd_outside_junta} we get that
    $\sum\limits_{S: S\not\subseteq \mathcal{J}}\widehat{f}(S)^2\leq \frac{\eps}{8}$.
    Define $h(x) = 1$ if $\sum\limits_{S: S\subseteq \mathcal{J}}\widehat{f}(S)\chi_S(x)\geq 1/2$ and $0$ otherwise.
    Clearly, $h$ is a $\mathcal{J}$-junta and we have that
    \[
    \norm{f-h}_1
    =
    \norm{f-h}_2^2
    \leq 4\left\|f-\sum\limits_{S: S\subseteq \mathcal{J}}\widehat{f}(S)\chi_S(x)\right\|_2^2
    =4\sum\limits_{S: S\not\subseteq \mathcal{J}}\widehat{f}(S)^2\leq \eps.
    \]
    Finally, note that as the sum of $\left(1-\frac{1}{3d}\right)$-noisy influences of $f$ is at most $O(d)$ by Fact~\ref{fact:sum_of_noisy}, we get that
    \[
    \card{\mathcal{J}}\leq O\left(\frac{d}{\delta}\right).
    \qedhere
    \]
  \end{proof}

\subsection{Proof of Claim~\ref{claim:super_fried_main}}\label{sec:missing_proofs}
For the proof of Claim~\ref{claim:super_fried_main} we need a version of the hypercontractive inequality with the noise operator, as follows:
\begin{lemma}\label{lemma:hypercontractivity_noise}[\cite{ODonnell}]
  Let $q>2$, and $\rho\leq \frac{1}{\sqrt{q-1}}$. Then for all $f\colon\{0,1\}^n\to\mathbb{R}$ we have
  $\norm{T_{\rho} f}_q\leq \norm{f}_2$.
\end{lemma}

We also need the following claim, asserting that noise only decreases Talagrand's boundary.
\begin{claim}\label{claim:noise_decreases_tal}
  For $f\colon \{0,1\}^n\to\mathbb{R}$ and any noise operator $S = T_{\rho_1}\otimes\ldots T_{\rho_n}$ we have
  $\Expect{x}{\norm{\nabla{S f}(x)}_p}\leq \Expect{x}{\norm{\nabla{f}(x)}_p}$.
\end{claim}
\begin{proof}
  \[\Expect{x}{\norm{\nabla (S f)(x)}_p}
  \leq
  \Expect{x}{\norm{S \nabla f(x)}_p}
  =\Expect{x}{\norm{\Expect{y\sim S x}{\nabla f(y)}}_p},
  \]
  and the result follows from the triangle inequality.
\end{proof}


  Take $(J,z)$ to be a random restriction so that each $j\in [n]$ is included in $J$ independently with probability $\frac{1}{3d}$,
  and consider $g = S_k'' f$.
  We calculate
  the contribution of $\bar{\mathcal{J}}$ to the level $1$ weight of a restriction of $g$:
  \begin{align*}
  d\Expect{J, z}{\sum\limits_{j\in J\cap \bar{\mathcal{J}}}\widehat{g_{\overline{J}\rightarrow z}}(\set{j})^2}
  =d\sum\limits_{j\in \overline{\mathcal{J}}}\Expect{J}{\sum\limits_{S}\widehat{g(S)}^2 1_{S\cap J = \{j\}}}
  &=d\sum\limits_{S}\frac{\card{S\cap \bar{\mathcal{J}}}}{3d}\left(1-\frac{1}{3d}\right)^{\card{S}-1}\widehat{g(S)}^2\\
  &=\frac{1}{3}\sum\limits_{S}\card{S\cap \bar{\mathcal{J}}}\left(1-\frac{1}{3d}\right)^{\card{S}-1}\widehat{g(S)}^2.
  \end{align*}
  Also,
  \begin{align*}
    \sum\limits_{j\not\in \mathcal{J}} \norm{S_k \partial_j f}_2^2
    =
    4\sum\limits_{S}
    \card{S\cap \bar{\mathcal{J}}}
    \left(1-\frac{1}{2^k}\right)^{-2}\rho_1^{2\card{S\cap \mathcal{J}}}
    \rho_2^{2\card{S\cap \bar{\mathcal{J}}}}\widehat{g(S)}^2,
  \end{align*}
  where $\rho_1 = \frac{1-1/d}{1-2/(3d)}$ and $\rho_2 = \frac{1-1/2^k}{1-1/2^k+1/(3d)}$. Using $\frac{a}{b}\leq \frac{a+c}{b+c}$ for $0\leq a\leq b$ and $c\geq 0$ we have
  \[
  \rho_1 \leq \frac{1-1/d+2/(3d)}{1-2/(3d)+2/(3d)} = 1-\frac{1}{3d},
  \qquad
  \rho_2 \leq \frac{1-1/2^k+(1/2^k-1/(3d))}{1-1/2^k+1/(3d)+(1/2^k-1/(3d))} = 1-\frac{1}{3d},
  \]
  and so
  \[
  \sum\limits_{j\not\in \mathcal{J}} \norm{S_k \partial_j f}_2^2\leq
  12\left(1-\frac{1}{d}\right)^{-2} d\Expect{J, z}{\sum\limits_{j\in J\cap \bar{\mathcal{J}}}\widehat{g_{\overline{J}\rightarrow z}}(\set{j})^2}
  \ll d\Expect{J, z}{\sum\limits_{j\in J\cap \bar{\mathcal{J}}}\widehat{g_{\overline{J}\rightarrow z}}(\set{j})^2}.
  \]
  Hence our
  goal of upper bounding the left hand side in Claim~\ref{claim:super_fried_main} reduces to
  upper bounding the contribution of $\mathcal{J}$ to the level $1$ of random restrictions of $g$. For $J, z$ define
 \[
  L_{J,z} = \sett{j\in J\cap \bar{\mathcal{J}}}{\card{\widehat{g_{\overline{J}\rightarrow z}}(\set{j})} \leq \tau},
  \qquad
  H_{J,z} = \sett{j\in J\cap \bar{\mathcal{J}}}{\card{\widehat{g_{\overline{J}\rightarrow z}}(\set{j})} > \tau}.
  \]
  Then
  \[
  \Expect{J, z}{\sum\limits_{j\in J\cap \bar{\mathcal{J}}}\widehat{g_{\overline{J}\rightarrow z}}(\set{j})^2}
  =\Expect{J, z}{\sum\limits_{j\in L_{J,z}}\widehat{g_{\overline{J}\rightarrow z}}(\set{j})^2}
  +\Expect{J, z}{\sum\limits_{j\in H_{J,z}}\widehat{g_{\overline{J}\rightarrow z}}(\set{j})^2}.
  \]

  \paragraph{Bounding the contribution from $L_{J,z}$.}
  For the first sum, we have
  \begin{align*}
  \sum\limits_{j\in L_{J,z}}\widehat{g_{\overline{J}\rightarrow z}}(\set{j})^2
  \leq
  \left(\sum\limits_{j\in L_{J,z}}\widehat{g_{\overline{J}\rightarrow z}}(\set{j})^2\right)^{p}
  &\leq \left(\sum\limits_{j\in L_{J,z}}\tau^{(2p-1)/p}\card{\widehat{g_{\overline{J}\rightarrow z}}(\set{j})}^{1/p}\right)^{p}\\
  &\leq \tau^{2p-1}\left(\sum\limits_{j}\card{\widehat{g_{\overline{J}\rightarrow z}}(\set{j})}^{1/p}\right)^{p}.
  \end{align*}
  In the last expression we have $\tau^{2p-1}$ times the $1/p$-norm of the vector $(\widehat{g_{\overline{J}\rightarrow z}}(\set{j}))_{j\in J}$.
  Thus, using the triangle inequality in a similar way to Lemma~\ref{lemma:tal_lvl1}, we may bound the last expression from above by
  \begin{equation}\label{eq3}
  \tau^{2p-1}\Expect{x}{\norm{\nabla{g_{\overline{J}\rightarrow z}}(x)}_{1/p}}.
  \end{equation}
  Taking expectation, we get that
  \[
  \Expect{J,z}{\sum\limits_{j\in L_{J,z}}\widehat{g_{\overline{J}\rightarrow z}}(\set{j})^2}
  \leq \tau^{2p-1}\Expect{J,z}{\Expect{x}{\norm{\nabla {g_{\overline{J}\rightarrow z}}(x)}_{1/p}}}
  = \tau^{2p-1}\Expect{x}{\Expect{J,z}{\norm{\nabla g_{\overline{J}\rightarrow z}(x)}_{1/p}}}.
  \]
  Fix $x,z$ and consider the vector $\nabla g(x,z)$.
  Note that
  \[
  \Expect{J}{\norm{\nabla g_{\overline{J}\rightarrow z}(x)}_{1/p}^{1/p}}
  =
  \Expect{J}{\sum\limits_{j}\card{\partial_j g(x,z)}^{1/p}1_{j\in J}}
  =\frac{1}{3d}\norm{\nabla g(x,z)}_{1/p}^{1/p},
  \]
  and it follows by Jensen's inequality that
  \[
  \Expect{J}{\norm{\nabla g_{\overline{J}\rightarrow z}(x)}_{1/p}}
  \leq
  \Expect{J}{\norm{\nabla g_{\overline{J}\rightarrow z}(x)}_{1/p}^{1/p}}^{p}
  \leq d^{-p} \norm{\nabla g(x,z)}_{1/p}.
  \]
  It follows that~\eqref{eq3} is upper bounded by $\tau^{2p-1}d^{-p}\Expect{x,z}{\norm{\nabla g(x,z)}_{1/p}}$.
  Using Claim~\ref{claim:noise_decreases_tal}, we have that this may be upper bounded by
  $\tau^{2p-1}d^{-p}\Expect{x,z}{\norm{\nabla f(x,z)}_{1/p}} = \tau^{2p-1}d^{-p} \Expect{x,z}{s_f(x,z)^p}$.



  \paragraph{Bounding the contribution from $H_{I,z}$.}
  We have
  \[
  \sum\limits_{j\in H_{J,z}}\widehat{g_{\overline{J}\rightarrow z}}(\set{j})^2
  \leq \tau^{-\frac{1}{100d}}\sum\limits_{j\in H_{J,z}}\card{\widehat{g_{\overline{J}\rightarrow z}}(\set{j})}^{2+\frac{1}{100d}}
  \leq \tau^{-\frac{1}{100d}}\sum\limits_{j\in J\cap\bar{\mathcal{J}}}\card{\widehat{g_{\overline{J}\rightarrow z}}(\set{j})}^{2+\frac{1}{100d}},
  \]
  and we next take expectation over $z$:
  \begin{equation}\label{eq5}
  \Expect{z}{\sum\limits_{j\in H_{J,z}}\widehat{g_{\overline{J}\rightarrow z}}(\set{j})^2}
  \leq
  \tau^{-\frac{1}{100d}}\sum\limits_{j\in J\cap\bar{\mathcal{J}}}\Expect{z}{\card{\widehat{g_{\overline{J}\rightarrow z}}(\set{j})}^{2+\frac{1}{100d}}}.
  \end{equation}
  Define the operator
  $T = T_{\rho_1}^{\mathcal{J}}T_{\rho_2}^{\overline{\mathcal{J}}}$
  where $\rho_1 = \frac{1-2/(3d)}{1-1/(3d)}$ and $\rho_2 = \frac{1-1/2^k+1/(3d)}{1-1/2^k+2/(3d)}$,
  and note that by the semi-group property of the noise operator (namely, the fact that $T_{\rho}T_{\rho'} = T_{\rho\rho'}$)
  we may write $g = T S_k' f$. Thus, $\widehat{g_{\overline{J}\rightarrow z}}(\set{j})$ is equal to
  \[
  \sum\limits_{S\subseteq \overline{J} \cup\{j\}, S\ni j} \widehat{g}(S)\chi_{S\setminus\{j\}}(z) =
  \sum\limits_{S\subseteq \overline{J} \cup\{j\}, S\ni j}
  \rho_1^{\card{S\cap \mathcal{J}}}
  \rho_2^{\card{S\cap \overline{\mathcal{J}}}}
  \widehat{S_k' f}(S)\chi_{S\setminus\{j\}}(z).
  \]
  Thus, using hypercontractivity, i.e.~Lemma~\ref{lemma:hypercontractivity_noise}, we have that
  \begin{align}\label{eq4}
  \Expect{z}{\card{\widehat{g_{\overline{J}\rightarrow z}}(\set{j})}^{2+\frac{1}{100d}}}
  &=\left\|\sum\limits_{S\subseteq \overline{J} \cup\{j\}, S\ni j}
  \rho_1^{\card{S\cap \mathcal{J}}}
  \rho_2^{\card{S\cap \overline{\mathcal{J}}}}
  \widehat{S_k' f}(S)\chi_{S\setminus\set{j}}(z)\right\|_{2+\frac{1}{100d}}^{2+\frac{1}{100d}}\notag\\
  &\leq \left\|\sum\limits_{S\subseteq \overline{J} \cup\{j\}, S\ni j}
  \rho_1^{\card{S\cap \mathcal{J}}}
  \rho_2^{\card{S\cap \overline{\mathcal{J}}}}
  \left(1-\frac{1}{3d}\right)^{-(\card{S}-1)}\widehat{S_k' f}(S)\chi_{S\setminus\set{j}}(z)\right\|_2^{2+\frac{1}{100d}}\notag\\
  &\leq \left\|\sum\limits_{S\ni j}\widehat{S_k' f}(S)\chi_{S\setminus\set{j}}(z)\right\|_2^{2+\frac{1}{100d}}.
  \end{align}
  In the last inequality, we used Parseval's equality and that fact that $\rho_1,\rho_2\leq 1-\frac{1}{3d}$.
  %
  %
%
%
%
%
%
%
%
%
%
%
  Note that
  \[
  \left\|\sum\limits_{S\ni j}\widehat{S_k' f}(S)\chi_{S\setminus\set{j}}(z)\right\|_2^{2+\frac{1}{100d}}
  \ll \norm{S_k' \partial_j f}_2^{2+\frac{1}{100d}},
  \]
  so plugging~\eqref{eq4} into~\eqref{eq5} yields that
  \[
  \Expect{z}{\sum\limits_{j\in H_{J,z}}\widehat{g_{\overline{J}\rightarrow z}}(\set{j})^2}\ll
  \tau^{-\frac{1}{100d}}\sum\limits_{j\in J\cap\bar{\mathcal{J}}}\norm{S_k' \partial_j f}_2^{2+\frac{1}{100d}}.
  \]
  Taking expectation over $J$ now gives that
  \[
  \Expect{z,J}{\sum\limits_{j\in H_{J,z}}\widehat{g_{\overline{J}\rightarrow z}}(\set{j})^2}
  \ll
  \tau^{-\frac{1}{100d}}\sum\limits_{j\in\bar{\mathcal{J}}} \norm{S_k' \partial_j f}_2^{2+\frac{1}{100d}} \Expect{J}{1_{j\in J}}
  = \frac{1}{3d}\tau^{-\frac{1}{100d}}\sum\limits_{j\in\bar{\mathcal{J}}} \norm{S_k' \partial_j f}_2^{2+\frac{1}{100d}}.
  \]
\qed
\begin{thebibliography}{KKO18}
\bibitem[Bec75]{Bec75}
William Beckner.
\newblock Inequalities in {F}ourier analysis.
\newblock {\em Annals of Mathematics}, 102:159--182, 1975.

\bibitem[Ber67]{Bernstein}
Arthur~J Bernstein.
\newblock Maximally connected arrays on the n-cube.
\newblock {\em SIAM Journal on Applied Mathematics}, 15(6):1485--1489, 1967.

\bibitem[Bon70]{Bonami}
Aline Bonami.
\newblock \'{E}tude des coefficients de {Fourier} des fonctions de $l^p(g)$.
\newblock {\em Annales de l'Institut Fourier}, 20(2):335--402, 1970.

\bibitem[Bou02]{Bourgain}
Jean Bourgain.
\newblock On the distribution of the fourier spectrum of boolean functions.
\newblock {\em Israel Journal of Mathematics}, 131(1):269--276, 2002.

\bibitem[Fri98]{Friedgut}
Ehud Friedgut.
\newblock Boolean functions with low average sensitivity depend on few
  coordinates.
\newblock {\em Combinatorica}, 18(1):27--35, 1998.

\bibitem[Gro75]{grossLogSob}
Leonard Gross.
\newblock Logarithmic {S}obolev inequalities.
\newblock {\em American Journal of Mathematics}, 97(4):1061--1083, 1975.

\bibitem[Har66]{Harper2}
Lawrence~H Harper.
\newblock Optimal numberings and isoperimetric problems on graphs.
\newblock {\em Journal of Combinatorial Theory}, 1(3):385--393, 1966.

\bibitem[Har76]{Hart}
Sergiu Hart.
\newblock A note on the edges of the n-cube.
\newblock {\em Discrete Mathematics}, 14(2):157--163, 1976.

\bibitem[KKK{\etalchar{+}}]{KKKMS}
Esty Kelman, Subhash Khot, Guy Kindler, Dor Minzer, and Muli Safra.
\newblock Theorems of {KKL}, {F}riedgut, and {T}alagrand via random
  restrictions and log-sobolev inequality.
\newblock In {\em {ITCS} 2021}.

\bibitem[KKL88]{KKL}
Jeff Kahn, Gil Kalai, and Nathan Linial.
\newblock The influence of variables on boolean functions (extended abstract).
\newblock In {\em 29th Annual Symposium on Foundations of Computer Science,
  White Plains, New York, USA, 24-26 October 1988}, pages 68--80, 1988.

\bibitem[KKO18]{kindler2018gaussian}
Guy Kindler, Naomi Kirshner, and Ryan O’Donnell.
\newblock Gaussian noise sensitivity and fourier tails.
\newblock {\em Israel Journal of Mathematics}, 225:71--109, 2018.

\bibitem[KT14]{kitsos2014inequalities}
Christos~P Kitsos and Thomas~L Toulias.
\newblock Inequalities for the {F}isher’s information measures.
\newblock {\em Handbook of Functional Equations: Functional Inequalities},
  pages 281--313, 2014.

\bibitem[Lin64]{Lindsey}
John~H Lindsey.
\newblock Assignment of numbers to vertices.
\newblock {\em The American Mathematical Monthly}, 71(5):508--516, 1964.

\bibitem[Mar74]{Margulis}
Grigorii~Aleksandrovich Margulis.
\newblock Probabilistic characteristics of graphs with large connectivity.
\newblock {\em Problemy peredachi informatsii}, 10(2):101--108, 1974.

\bibitem[O'D14]{ODonnell}
Ryan O'Donnell.
\newblock {\em Analysis of boolean functions}.
\newblock Cambridge University Press, 2014.

\bibitem[Tal93]{Tal1}
Michel Talagrand.
\newblock Isoperimetry, logarithmic {S}obolev inequalities on the discrete
  cube, and {M}argulis' graph connectivity theorem.
\newblock {\em Geometric \& Functional Analysis GAFA}, 3(3):295--314, 1993.

\bibitem[Tal94]{Tal4}
Michel Talagrand.
\newblock On {R}usso's approximate zero-one law.
\newblock {\em The Annals of Probability}, 22(3):1576--1587, 1994.

\bibitem[Tal97]{Tal2}
Michel Talagrand.
\newblock On boundaries and influences.
\newblock {\em Combinatorica}, 17(2):275--285, 1997.

\end{thebibliography}\end{document}
